%==============================================================================
%  Statistical Rethinking -- Chapter 13
%  Multilevel Models: Varying Intercepts and Partial Pooling
%  Worked example on the Inside Airbnb "listings" data
%
%  Compile with:  pdflatex essay.tex   (twice, for references)
%==============================================================================
\documentclass[11pt,a4paper]{article}

%------------------------------- packages ------------------------------------
\usepackage[T1]{fontenc}
\usepackage[utf8]{inputenc}
\usepackage[margin=1in]{geometry}
\usepackage{amsmath,amssymb}
\usepackage{graphicx}
\usepackage{booktabs}
\usepackage{listings}
\usepackage{xcolor}
\usepackage{caption}
\usepackage[hidelinks]{hyperref}

%---------------------------- R code formatting ------------------------------
\definecolor{codegray}{gray}{0.95}
\definecolor{codegreen}{rgb}{0.0,0.45,0.15}
\definecolor{codeblue}{rgb}{0.10,0.20,0.65}

\lstdefinestyle{Rstyle}{
  language        = R,
  backgroundcolor = \color{codegray},
  basicstyle      = \ttfamily\footnotesize,
  keywordstyle    = \color{codeblue}\bfseries,
  commentstyle    = \color{codegreen}\itshape,
  stringstyle     = \color{purple},
  numbers         = left,
  numberstyle     = \tiny\color{gray},
  numbersep       = 8pt,
  breaklines      = true,
  frame           = single,
  framesep        = 4pt,
  rulecolor       = \color{gray},
  showstringspaces= false,
  tabsize         = 2,
  captionpos      = b,
}
\lstset{style=Rstyle}

%-------------------------------- metadata -----------------------------------
\title{\textbf{Partial Pooling Across Hosts:\\ A Multilevel Model for Airbnb Listing Prices}\\
       \large Statistical Rethinking, Chapter 13 --- Worked Exercise}
\author{TalRamot46\thanks{Replace with your full name, ID, and course details.}}
\date{\today}

\begin{document}
\maketitle

\begin{abstract}
% TEMPLATE: 120--200 words. Say (i) what question you ask, (ii) what data you use,
% (iii) which model you fit, (iv) the headline result, (v) why partial pooling matters here.
This essay applies the varying-intercepts multilevel model of Chapter~13 of
\emph{Statistical Rethinking} to \textcolor{red}{[data set]}. The outcome of
interest is \textcolor{red}{[outcome]}, grouped by \textcolor{red}{[cluster
variable]}. We fit \textcolor{red}{[model]} using Hamiltonian Monte Carlo and
compare the resulting cluster-level estimates to their no-pooling and
complete-pooling counterparts. We find \textcolor{red}{[headline result]}, which
illustrates \textcolor{red}{[the shrinkage / regularisation lesson]}.
\end{abstract}

\tableofcontents
\bigskip

%==============================================================================
\section{Introduction and Research Question}
%==============================================================================
% TEMPLATE: 3--5 paragraphs.
%   (a) Motivate the substantive question in plain language.
%   (b) State why the data are *clustered* -- i.e. why observations are not
%       independent -- since this is the whole justification for Chapter 13.
%   (c) State the research question precisely, as a statement about a parameter.
%   (d) Preview the structure of the essay.

Chapter~13 of \emph{Statistical Rethinking} introduces multilevel (hierarchical)
models as a principled compromise between two extremes: \emph{complete pooling},
which ignores cluster identity entirely, and \emph{no pooling}, which estimates
each cluster in isolation. \textcolor{red}{[Extend this motivation to your own
setting.]}

The data used here are \textcolor{red}{[describe source]}. Each row is a single
listing, but listings are nested within hosts: a single host may operate many
listings, and those listings plausibly share a common pricing strategy. Treating
the rows as independent would therefore overstate the effective sample size and
understate uncertainty.

Our research question is:
\begin{quote}
\textcolor{red}{[State it as a question about $\bar{\alpha}$, $\sigma_{\text{host}}$,
or the ratio between them.]}
\end{quote}

The remainder of the essay proceeds as follows. Section~\ref{sec:data} describes
the data and preprocessing. Section~\ref{sec:model} states the statistical model
and justifies each prior. Section~\ref{sec:fit} reports the fitting procedure and
diagnostics. Section~\ref{sec:results} presents the posterior. Section~\ref{sec:pooling}
examines shrinkage directly, and Section~\ref{sec:conclusion} concludes.

%==============================================================================
\section{Data}\label{sec:data}
%==============================================================================

\subsection{Source and Variables}
% TEMPLATE: Name the source, the scrape date, the number of raw rows, and the
% number of rows surviving each filter. State units for every variable.

The raw file \texttt{listings.csv} contains
\textcolor{red}{[$N_{\text{raw}}$]} listings and 91 columns. The variables used
in this analysis are summarised in Table~\ref{tab:variables}.

\begin{table}[htbp]
\centering
\caption{Variables used in the analysis.}
\label{tab:variables}
\begin{tabular}{@{}llp{6.5cm}@{}}
\toprule
\textbf{Variable} & \textbf{Type} & \textbf{Description} \\
\midrule
\texttt{host\_id}              & integer & Identifier of the host; defines the cluster. \\
\texttt{price}                 & string  & Nightly price as scraped, e.g.\ \verb|$120.00|. \\
\texttt{review\_scores\_rating}& real    & Mean guest rating on a 0--5 scale. \\
\texttt{number\_of\_reviews}   & integer & Count of reviews; used as a filter. \\
\midrule
\texttt{price\_clean}          & real    & \emph{Derived.} \texttt{price} parsed to a number. \\
\texttt{log\_price}            & real    & \emph{Derived.} $\log(\texttt{price\_clean})$; the outcome. \\
\texttt{host\_idx}             & integer & \emph{Derived.} Contiguous index $1,\dots,J$ over hosts. \\
\bottomrule
\end{tabular}
\end{table}

\subsection{Preprocessing}
% TEMPLATE: Justify EACH decision below -- a grader will look for the reasoning,
% not just the code.
%   * Why filter number_of_reviews > 0?
%   * Why log-transform price? (skew; multiplicative effects; support on R)
%   * Why subsample to n = 2000? (compute cost -- and what it costs you)
%   * How many hosts J remain, and what is the distribution of listings per host?
%     This last point is CRITICAL: partial pooling only bites when some clusters
%     are small.

Three preprocessing choices deserve comment. First, listings with no reviews are
dropped, since \textcolor{red}{[reason]}. Second, price is modelled on the log
scale because \textcolor{red}{[reason]}; this also guarantees that the Gaussian
likelihood places no mass on negative prices once exponentiated. Third, the data
are subsampled to $n = 2000$ rows to keep sampling times manageable, at the cost
of \textcolor{red}{[reason]}.

\textbf{Cluster structure.} After filtering, there are $J = \textcolor{red}{[J]}$
distinct hosts across $n = 2000$ listings, a mean of
\textcolor{red}{[$n/J$]} listings per host. \textcolor{red}{[Report the
distribution --- how many hosts contribute exactly one listing? Those are the
clusters that shrink the most.]}

\begin{figure}[htbp]
\centering
% \includegraphics[width=0.8\textwidth]{figures/listings_per_host.pdf}
\fbox{\parbox[c][4.5cm][c]{0.8\textwidth}{\centering \textcolor{gray}{%
[Figure placeholder: histogram of listings per host, and histogram of
\texttt{log\_price}.]}}}
\caption{Distribution of listings per host (left) and of \texttt{log\_price} (right).}
\label{fig:eda}
\end{figure}

%==============================================================================
\section{The Model}\label{sec:model}
%==============================================================================

\subsection{Statistical Specification}

Let $y_i$ denote \texttt{log\_price} for listing $i = 1,\dots,n$, and let
$j[i] \in \{1,\dots,J\}$ denote the host of listing $i$. The varying-intercepts
model is
\begin{align}
  y_i          &\sim \mathrm{Normal}(\mu_i,\ \sigma) \\
  \mu_i        &= \alpha_{j[i]} \\
  \alpha_j     &\sim \mathrm{Normal}(\bar{\alpha},\ \sigma_{\text{host}})
                  &&\text{[adaptive prior]} \\
  \bar{\alpha} &\sim \mathrm{Normal}(4.5,\ 1) &&\text{[hyper-prior]} \\
  \sigma_{\text{host}} &\sim \mathrm{Exponential}(1) &&\text{[hyper-prior]} \\
  \sigma       &\sim \mathrm{Exponential}(1).
\end{align}

The third line is what makes this a multilevel model. The prior for each host
intercept $\alpha_j$ is not fixed in advance but is itself estimated from the
data through $\bar{\alpha}$ and $\sigma_{\text{host}}$. Hosts therefore
\emph{learn from one another}: a host with a single listing is pulled towards the
population mean $\bar{\alpha}$, whereas a host with many listings largely retains
its own estimate.

\subsection{Prior Justification}
% TEMPLATE: Justify each prior on the scale of the outcome. Do NOT just say
% "weakly informative" -- show what it implies.
%   * a_bar ~ Normal(4.5, 1): on the log scale, exp(4.5) is about 90 currency
%     units per night, and +/- 2 SD spans exp(2.5) to exp(6.5).  Is that a
%     plausible price range for this city?  Say so.
%   * sigma_host ~ Exponential(1): mean 1 on the log scale -- i.e. hosts are
%     expected to differ by a factor of about e.  Is that generous or tight?
%   * sigma ~ Exponential(1): residual variation *within* a host.
%   * Optionally run a prior predictive simulation and show it as a figure.

\begin{figure}[htbp]
\centering
% \includegraphics[width=0.7\textwidth]{figures/prior_predictive.pdf}
\fbox{\parbox[c][4cm][c]{0.7\textwidth}{\centering \textcolor{gray}{%
[Figure placeholder: prior predictive simulation of prices implied by the priors
above, on the natural (currency) scale.]}}}
\caption{Prior predictive check. Priors should permit plausible prices without
concentrating mass on absurd ones.}
\label{fig:prior-pred}
\end{figure}

\subsection{Comparison Models}
% TEMPLATE: Chapter 13's central argument is comparative. Fit at least the two
% reference models below so you have something to compare shrinkage against.
\begin{itemize}
  \item \textbf{Complete pooling:} $\mu_i = \alpha$, one intercept for all listings.
  \item \textbf{No pooling:} $\mu_i = \alpha_{j[i]}$ with $\alpha_j \sim \mathrm{Normal}(4.5, 1)$
        fixed --- each host estimated independently.
  \item \textbf{Partial pooling:} the model above.
\end{itemize}
Compare them with \texttt{compare(m1, m2, m3, func=WAIC)} and discuss the
effective number of parameters, which is the clearest quantitative signature of
regularisation.

%==============================================================================
\section{Fitting and Diagnostics}\label{sec:fit}
%==============================================================================

The model was fit with \texttt{ulam()}, which compiles the formula to Stan and
samples via Hamiltonian Monte Carlo: 4 chains, 2000 iterations each (half warm-up).
A \texttt{brms} implementation of the same structure is included for comparison.

\subsection{Diagnostics}
% TEMPLATE: Report and interpret, do not merely list.
%   * R-hat for all parameters (target: <= 1.01)
%   * n_eff / ESS -- flag any parameter below ~200
%   * Divergent transitions -- if present, discuss non-centred parameterisation
%     (Chapter 13's "Meta" section), and show the reparameterised code.
%   * Trace plots / trank plots for a_bar, sigma_host, sigma.

\begin{table}[htbp]
\centering
\caption{Convergence diagnostics for the population-level parameters.}
\label{tab:diagnostics}
\begin{tabular}{@{}lrrrr@{}}
\toprule
\textbf{Parameter} & \textbf{Mean} & \textbf{SD} & \textbf{ESS} & $\hat{R}$ \\
\midrule
$\bar{\alpha}$            & \textcolor{red}{--} & \textcolor{red}{--} & \textcolor{red}{--} & \textcolor{red}{--} \\
$\sigma_{\text{host}}$    & \textcolor{red}{--} & \textcolor{red}{--} & \textcolor{red}{--} & \textcolor{red}{--} \\
$\sigma$                  & \textcolor{red}{--} & \textcolor{red}{--} & \textcolor{red}{--} & \textcolor{red}{--} \\
\bottomrule
\end{tabular}
\end{table}

%==============================================================================
\section{Results}\label{sec:results}
%==============================================================================

\subsection{Population-Level Parameters}
% TEMPLATE: Interpret on the natural scale. exp(a_bar) is the typical nightly
% price of a typical host. Give a 89% compatibility interval (the book's default)
% and say what it means.

\subsection{Variance Decomposition}
% TEMPLATE: The most informative single number in this model is
%   ICC = sigma_host^2 / (sigma_host^2 + sigma^2),
% the share of variation in log-price attributable to *between-host* differences.
% Compute it from the posterior draws (so you get an interval, not a point) and
% interpret it substantively.
\begin{equation}
  \mathrm{ICC} \;=\; \frac{\sigma_{\text{host}}^{2}}
                          {\sigma_{\text{host}}^{2} + \sigma^{2}}
  \;=\; \textcolor{red}{[\,\cdot\,]}
  \quad (89\%\ \text{CI}: \textcolor{red}{[\,\cdot\,,\,\cdot\,]}).
\end{equation}

\subsection{Host-Level Intercepts}
\begin{figure}[htbp]
\centering
% \includegraphics[width=0.85\textwidth]{figures/caterpillar.pdf}
\fbox{\parbox[c][5cm][c]{0.85\textwidth}{\centering \textcolor{gray}{%
[Figure placeholder: caterpillar plot of the $\alpha_j$ posteriors, ordered by
mean, with $\bar{\alpha}$ as a horizontal reference line.]}}}
\caption{Posterior means and 89\% intervals for host intercepts $\alpha_j$.}
\label{fig:caterpillar}
\end{figure}

%==============================================================================
\section{Shrinkage and Partial Pooling}\label{sec:pooling}
%==============================================================================
% TEMPLATE: This is the heart of a Chapter 13 essay. Reproduce the book's
% Figure 13.1-style plot: raw (no-pooling) host means on the x-axis against
% partially pooled posterior means on the y-axis, with point size proportional
% to the number of listings for that host.
%
% Then make three claims and support each with the plot:
%   1. Shrinkage is towards a_bar, not towards zero.
%   2. Shrinkage is strongest for hosts with few listings.
%   3. Shrinkage is strongest for hosts whose raw mean is far from a_bar.
%
% Close by connecting this to out-of-sample prediction: the pooled estimates are
% *better* predictors of a host's next listing, even though they fit the observed
% data worse. That trade-off is the chapter's central lesson.

\begin{figure}[htbp]
\centering
% \includegraphics[width=0.75\textwidth]{figures/shrinkage.pdf}
\fbox{\parbox[c][6cm][c]{0.75\textwidth}{\centering \textcolor{gray}{%
[Figure placeholder: no-pooling estimate vs.\ partially pooled estimate per host;
point size $\propto$ listings per host; 45$^\circ$ line and $\bar{\alpha}$ shown.]}}}
\caption{Shrinkage of host-level estimates towards the population mean.}
\label{fig:shrinkage}
\end{figure}

%==============================================================================
\section{Limitations}
%==============================================================================
% TEMPLATE: Be specific and honest. Candidates:
%   * The model has no predictors -- it is an intercept-only decomposition of
%     variance, so "host effects" absorb room type, neighbourhood, capacity, etc.
%     A varying-intercepts model with these covariates would separate them.
%   * Subsampling to 2000 rows fragments the cluster structure: hosts with many
%     listings in the full data may appear only once in the sample.
%   * Prices are scraped snapshots, not transacted prices.
%   * Gaussian residuals on the log scale still assume symmetry; check with a
%     posterior predictive plot.

%==============================================================================
\section{Conclusion}\label{sec:conclusion}
%==============================================================================
% TEMPLATE: 2--3 paragraphs. Answer the research question stated in Section 1,
% in the same terms. Then state what Chapter 13's machinery bought you that a
% single-level regression would not have.

%==============================================================================
\appendix
\section{R Code}
%==============================================================================

\subsection{Multilevel Model via \texttt{rethinking::ulam}}
\lstinputlisting[caption={\texttt{rethinking\_code.R} --- varying intercepts on
\texttt{log\_price} by host.}, label={lst:ulam}]{rethinking_code.R}

\subsection{Equivalent Model via \texttt{brms}}
\lstinputlisting[caption={\texttt{code.R} --- the same varying-intercepts
structure in \texttt{brms} syntax.}, label={lst:brms}]{code.R}

\subsection{Reproducibility}
% TEMPLATE: Paste sessionInfo() output, the RNG seed, and Stan/cmdstan versions.
% Note that slice_sample() without set.seed() makes the analysis irreproducible --
% add set.seed() to the scripts before submitting.

%==============================================================================
\begin{thebibliography}{9}
%==============================================================================
\bibitem{mcelreath}
McElreath, R. (2020). \emph{Statistical Rethinking: A Bayesian Course with
Examples in R and Stan}, 2nd ed. CRC Press. Chapter 13.

\bibitem{burkner}
B\"urkner, P.-C. (2017). brms: An R Package for Bayesian Multilevel Models Using
Stan. \emph{Journal of Statistical Software}, 80(1), 1--28.

\bibitem{stan}
Stan Development Team (2024). \emph{Stan Modeling Language Users Guide and
Reference Manual}.

\bibitem{insideairbnb}
Inside Airbnb. \url{http://insideairbnb.com/get-the-data/}. Accessed
\textcolor{red}{[date]}.
\end{thebibliography}

\end{document}
